% sections/travail_realise.tex
% Phase 10: Travail realise (RAPP-05)

\section{Travail r\'ealis\'e}
\label{sec:travail_realise}

Le projet s'est d\'eroul\'e en cinq phases~: collecte des donn\'ees,
construction du graphe, analyses par question de recherche, r\'edaction
du rapport, et d\'eveloppement d'un prototype web interactif compl\'ementaire.

\paragraph{Collecte et pr\'eparation des donn\'ees.}
Les 401\,709 transferts WETH du 5 ao\^ut 2024 ont \'et\'e extraits de la
table publique \texttt{crypto\_polygon.transactions} sur Google
BigQuery~\cite{bigquery_public}. Le r\'esultat, stock\'e dans le fichier
CSV \texttt{polygon\_weth\_crash\_2024-08-05.csv} (74\,Mo), a \'et\'e soumis
\`a un pr\'etraitement d'agr\'egation~: les ar\^etes multiples entre une m\^eme
paire d'adresses (source, target) sont fusionn\'ees par somme des poids
(volume WETH) et comptage des transactions. Un filtrage
\textit{dust} \`a $10^{-12}$~WETH supprime les micro-transferts n\'egligeables.
Le graphe r\'esultant comprend 12\,880 sommets et 39\,999 ar\^etes orient\'ees
pond\'er\'ees.

\paragraph{Construction du graphe.}
Le module \texttt{graph\_builder.py} sert d'interface entre les donn\'ees
brutes et les six analyses. Sa fonction \texttt{load\_and\_build()} retourne
un triplet \texttt{(df, digraph, graph)}~: le DataFrame Pandas, un graphe
orient\'e pond\'er\'e (\texttt{DiGraph}), et un graphe non orient\'e
(\texttt{Graph}) pour RQ3 et RQ5. Toutes les analyses travaillent ainsi sur
le m\^eme jeu de donn\'ees.

\paragraph{Analyses par question de recherche.}
Six modules ind\'ependants (\texttt{src/rq1\_*.py} \`a \texttt{src/rq6\_*.py})
r\'ealisent chacun une analyse et produisent deux figures PNG \`a 200~DPI
(16~au total). Le module \texttt{src/style.py} assure la coh\'erence
graphique~: fond sombre pour les graphes de r\'eseau, fond clair pour les
graphiques analytiques. Le script \texttt{run\_all.py} ex\'ecute les six
analyses \`a la suite.

\paragraph{R\'edaction du rapport.}
Le rapport est r\'edig\'e en \LaTeX{} et compil\'e avec pdf\LaTeX{} + Biber.
Le fichier principal \texttt{main.tex} inclut 17~fichiers \texttt{.tex}
(sections et annexes). La bibliographie est scind\'ee en deux~: r\'ef\'erences
acad\'emiques et webographie, filtr\'ees par mot-cl\'e via \texttt{biblatex}.

\paragraph{Prototype web interactif.}
En parall\`ele du rapport, on a construit un prototype web pour explorer
le r\'eseau sans installer Python. Il est d\'eploy\'e \`a l'adresse
\url{https://crypto-graph-explorer.vercel.app}~\cite{cryptographexplorer_web}.
La visualisation force-directed utilise D3.js, l'interface est faite en
React + Vite avec les composants shadcn/ui. Le prototype lit un CSV de
m\'etriques pr\'e-calcul\'ees (\texttt{nodes\_metrics.csv}, $12\,880$
n\oe{}uds, degr\'e, PageRank et num\'ero de communaut\'e), et un panneau
lat\'eral donne acc\`es aux six RQ et aux cinq analyses compl\'ementaires
de la section~\ref{sec:complementaires}. Les chiffres exacts ($G$,
$\sigma$, $Q$) viennent uniquement de l'ex\'ecution Python sur les
donn\'ees brutes~; le prototype, lui, calcule des versions approch\'ees
adapt\'ees au navigateur. L'annexe~\ref{ann:manuel} d\'ecrit l'acc\`es
et ses limites plus en d\'etail.

\begin{table}[H]
\centering
\caption{R\'epartition des t\^aches du projet}
\label{tab:repartition}
\begin{tabular}{lll}
\toprule
\textbf{Membre} & \textbf{T\^ache} & \textbf{Contribution} \\
\midrule
Thibaud THOMAS-LAMOTTE & Collecte BigQuery, RQ1, RQ2, RQ3 & Topologie, centralit\'e, communaut\'es \\
Youcef HEMISSI & RQ4, RQ5, RQ6, rapport \LaTeX{} & Temporel, petit-monde, flux \\
\bottomrule
\end{tabular}
\end{table}
